\section{Compatible piecewise-smooth tests and critical tangent pairings}
\label{sec:r13-compatible-ps-tangent-pairing}

The currents in the first part of this section are real
supercurrents on an open subset $\Omega\subset\mathbf R^N$,
with $N\ge2$.  They are symmetric, middle-positive,
separately $d'$- and $d''$-closed, and have bidimension
$(2,2)$.  Write
\[
             \beta=\sum_i d'x_i\wedge d''x_i,
                      \mu_T=\tfrac12T\wedge\beta^2.
 \tag{PT1}
\]
All coefficient measures are controlled by this trace, as in
MT4.  No weak/strong positivity identification is made.

\paragraph{Nonduplication status.}
The uniform $O(\sqrt\epsilon)$ piecewise-test approximation
mechanism is retained: R6,
\texttt{lem:r6-weighted-pws-error}, records it as a direct
corollary of R4 Lemmas 71.1--71.2 and Theorem 71.4, even
with bounded Borel multipliers.  Its real bidimension-two
proof is written below to verify precisely the tangent-limit
interface used here.  The nodal construction likewise makes
the retained Courant scalar generation explicit in the
homogeneous constant-tensor case.  These mechanisms are not
counted as new principal theorems.  The new composition is
their use with MT and LE to give the complete realized
critical-tangent pairing in PT11.

A \emph{compatible finite-PS test} is a compactly supported
ambient $(2,2)$ form on a finite polyhedral complex covering
a neighborhood of its support, with smooth extensions near
the closed cells, whose tangential pullbacks agree on every
common face.  This is a supplied genuine ambient presentation,
not independent coefficient data attached to different source
labels.  Bounded first derivatives on a slightly larger
compact neighborhood follow from the finite closed-cell data.

\subsection{Uniform control of normal factors near a wall}

Let $H=\{\ell=0\}$ be an affine hyperplane, with
$|d\ell|=1$, and put
$n'=d'\ell$, $n''=d''\ell$.  A $(2,2)$ coefficient
tensor belongs to its \emph{normal-factor subspace} if it
lies in
\[
 I_H=n'\wedge\Lambda^{1,2}
                     +n''\wedge\Lambda^{2,1}.
 \tag{PT2}
\]
This is exactly the kernel of tangential restriction to $H$.

\begin{lemma}[Normal strip bounds]
\label{lem:r13-normal-strip-bounds}
Fix compact sets $K\subset\operatorname{int}K'\Subset\Omega$.
For every such $T$ and every sufficiently small $\epsilon>0$,
\[
 \begin{split}
 (T\wedge n'\wedge n''\wedge\beta)
       (K\cap\{ |\ell|\le\epsilon\})
             &\le C\epsilon\,\mu_T(K'),\\
 |\langle T,\psi\rangle|
             &\le C\|\psi\|_\infty\sqrt\epsilon\,
                              \mu_T(K')
 \end{split}
 \tag{PT3}
\]
whenever the bounded Borel test $\psi$ is supported in that
strip and takes values in $I_H$.  The constants depend on
the fixed geometry and cutoff, not on $T$.  The coefficient
measures involving at least one normal two-vector index have
zero restriction to $H$ itself.  The trace $\mu_T$ may
nevertheless charge $H$.
\end{lemma}

\begin{proof}
Choose $0\le\chi\in C_c^\infty(\operatorname{int}K')$
equal to one near $K$.  Choose smooth convex functions
$u_\epsilon$ on the real line such that
\[
 u_\epsilon''(s)\ge\epsilon^{-1}
          \quad (|s|\le\epsilon),\qquad
 |u_\epsilon(s)|\le C(1+|s|)
 \tag{PT4}
\]
on the fixed relevant range, with $C$ independent of
$\epsilon$.  For instance rescale one convex smoothing
of $|s|$ whose second derivative has compact support and
is bounded below on the unit interval.  Then
$D(u_\epsilon\circ\ell)=u_\epsilon''(\ell)n'\wedge n''$.
Positivity and closedness give
\[
 \begin{split}
 \epsilon^{-1}
   (T\wedge n'\wedge n''\wedge\beta)
          (K\cap\{|\ell|\le\epsilon\})
 &\le\langle T,\chi D(u_\epsilon\circ\ell)\wedge\beta\rangle\\
 &=\langle T,(u_\epsilon\circ\ell)D\chi\wedge\beta\rangle\\
 &\le C\mu_T(K').
 \end{split}
 \tag{PT5}
\]
The last inequality follows from the bounded scalar function
in PT4 and the fixed middle-order bound for $D\chi\wedge\beta$.
This is a distributional Stokes identity with smooth functions;
no current is restricted to the wall in deriving it.

Use an orthonormal frame with first covector $d\ell$.
The cocoefficient matrix of $T$ is positive semidefinite on
two-vectors.  Its diagonal entries whose index contains this
first covector are bounded by the normal trace in PT5; its
other diagonal entries are bounded by the total trace.
For a Borel set $B$, positivity and Radon polarization give
\[
           |\mu_{IJ}|(B)
                   \le\sqrt{\mu_{II}(B)\mu_{JJ}(B)}.
 \tag{PT6}
\]
Each coefficient of an $I_H$-valued tensor has such a normal
index in at least one column.  Apply PT5 to that diagonal
entry and the trace bound to the other.  The finite sum of
coefficient estimates proves the second inequality of PT3.
Letting $\epsilon\downarrow0$ proves the assertion on $H$.
It only removes normal-factor coefficient measures there,
not the tangential trace mass.
\end{proof}

\subsection{A finite-wall approximation with all lower faces retained}

\begin{lemma}[Uniform smooth approximation of compatible tests]
\label{lem:r13-uniform-compatible-ps-approximation}
Let $\theta$ be a compatible finite-PS compact test.  There
are compactly supported smooth coefficient convolutions
$\theta_\epsilon$ in one fixed compact neighborhood and a
constant $C_\theta$ such that
\[
       |\langle T,\theta-\theta_\epsilon\rangle|
                 \le C_\theta\sqrt\epsilon\,\mu_T(K')
 \tag{PT7}
\]
for every current in the stated class.  The pairing with
$\theta$ is independent of ambient representatives of its
lower-face values.
\end{lemma}

\begin{proof}
Take the finite list of affine hyperplanes containing the
coefficient facets, and refine by their arrangement.  Fix
one smooth closed-cell extension for each resulting full
cell.  At a point of a lower face, choose an incident full
cell to give a bounded Borel representative.  Any two such
choices differ by a sum of tensors in $I_H$ for hyperplanes
containing that face: adjacent incident cells have this
property by tangential compatibility, and the local full
cell adjacency through facets is connected.  Alternatively
the kernel of restriction to the face is the sum of the
ideals generated by its conormals.  PT3 makes all these
differences pair to zero on the face.  The chosen Borel
representative therefore gives the canonical pairing.

Extend by zero outside a fixed compact neighborhood of the
support, where the original test already vanishes.  Let
$\theta_\epsilon$ be its coefficient convolution with a
smooth kernel supported in $|a|\le\epsilon$.
For a generic translation vector $a$, traverse the segment
from $x$ to $x-a$.  Within cells, its total smooth variation
is at most $C|a|$.  At a crossed facet $H$ the jump is in
$I_H$, because its tangential restriction is zero.  All
these jump coefficients are uniformly bounded by the finite
closed-cell bounds.  The hyperplane crossed by the segment
has distance at most $|a|$ from $x$.  A generic segment
crosses at most one facet of each of the finitely many
hyperplanes; crossings at lower faces can be omitted from
the translation integral, since they form a Lebesgue-null
set of translations.  If $x$ itself lies on lower faces,
the initial choice of incident cell is handled by the same
finite incident-cell chain just used above.

In a fixed coordinate basis this proves the decomposition
\[
       \theta_\epsilon-\theta
             =R_\epsilon+\sum_H J_{H,\epsilon},
 \tag{PT8}
\]
where $\|R_\epsilon\|_\infty\le C\epsilon$, each
$J_{H,\epsilon}$ is a bounded Borel $I_H$-valued tensor
with norm at most $C$, and its support lies in a fixed
multiple of the $\epsilon$-strip around $H$.  Indeed the
conormal of an affine hyperplane is constant, so transporting
a jump from its crossing point to $x$ preserves membership
in $I_H$.  A fixed linear splitting of this finite-dimensional
subspace gives bounded coefficient representatives before
integrating over the translations.  There are finitely many
walls and no accumulating crossing sequence.

The ordinary coefficient trace bound controls the first term
in PT8 by $C\epsilon\mu_T(K')$.  PT3 controls each remaining
term by $C\sqrt\epsilon\mu_T(K')$.  Their finite sum proves
PT7, including the strata where several coefficient walls
meet.  No simultaneous covector primitive is used or needed.
\end{proof}

\begin{theorem}[Weak continuity on compatible finite-PS tests]
\label{thm:r13-compatible-ps-weak-continuity}
Suppose $T_j\to T_0$ on all compact smooth real tests and
the currents satisfy the stated middle positivity and both
closedness equations, with uniformly bounded trace on each
compact subset.  Then for every fixed compatible finite-PS
compact test,
\[
                     \langle T_j,\theta\rangle
                              \longrightarrow
                     \langle T_0,\theta\rangle.
 \tag{PT9}
\]
The limit may charge coefficient walls.  No zero-wall-mass
hypothesis is required.
\end{theorem}

\begin{proof}
The limit has the same middle positivity and closedness and
the same local trace bounds.  For fixed $\epsilon$ the
convoluted test is smooth, so its pairings converge.  PT7
bounds the two approximation errors uniformly in $j$.
First pass to the limit in $j$, then let $\epsilon$ tend
to zero.  This proves PT9.
\end{proof}

\subsection{A constructive ambient extension for homogeneous tensors}

\begin{lemma}[Nodal extension from a geometric subfan]
\label{lem:r13-nodal-homogeneous-form-extension}
Let $\Sigma$ be a finite rational fan in $\mathbf R^N$.
Suppose a constant $(p,q)$ tensor is given on the span of
each cone, with identical tangential restrictions on common
faces.  Then it has a genuine ambient homogeneous
coefficient-degree-zero finite-PS extension.  If the data
are column-symmetric, the extension may be chosen symmetric.
This lemma extends regular forms, not their primitives.
\end{lemma}

\begin{proof}
Refine a complete rational hyperplane fan containing all
equations for the spans of the supplied cones, extensions
of their relative facet equations, and the coordinate
hyperplanes, and make a
common simplicial refinement.  This gives a complete
simplicial fan $\Delta$ such that the supplied support is
a subfan after subdivision.  The coordinate hyperplanes
remove lineality before triangulation.  Choose a vector
$v_r$ on every ray and its continuous homogeneous PL
nodal function $\varphi_r$: on a cone containing that ray
it is the corresponding coordinate, and on other cones
it is zero.  These formulas agree on shared faces.

Fix a global order of the rays.  For strictly increasing
ray multi-indices $I,J$ of lengths $p,q$ in that order,
consider their union.  If its rays span a face $\tau$
of $\Delta$ lying in the supplied subfan, let $a_{IJ}$
be the value of the supplied tensor on the corresponding
ordered ray vectors in its two columns.  The value is
independent of a containing supplied cone by compatibility.
If their union is not contained in any cone, or its smallest
face is not supplied, set $a_{IJ}=0$.  The finite sum
\[
     \widehat E=\sum_{I,J}a_{IJ}
                  d'_P\varphi_I\wedge d''_P\varphi_J
 \tag{PT10}
\]
with the standard column sign convention defines an actual
ambient finite-PS form.  On any full simplicial cone the
nonzero nodal differentials are its dual ray basis.  PT10
therefore has exactly the assigned tensor coefficients.
On a supplied cone every relevant smaller ray face is
supplied, so its restriction is the original tensor.
The same argument proves all shared-face compatibilities.
Column symmetry is preserved by making the symmetric
coefficient assignments.  This explicit scalar construction
does not presume an extension of independently chosen
source primitives or identify coincident source labels.
\end{proof}

Only geometric single-valued face data are covered by this
extension lemma.  Different actual source labels with
incompatible values on the same geometric image do not
define its input.  LE's leading energy already has a
supplied ambient presentation by the scalar first jets;
PT10 is an additional constructive extension option.

\subsection{The central scalar as a tangent-current pairing}

Return to the protected original setting of MT and LE.
Let $T$ be the full-fiber real background transpose, and
let $T_{\delta_j}=\delta_j^2(a_{\delta_j})_*T$ converge
to any critical tangent $T_0$, where
$a_\delta(x)=x/\delta$.  Let $E$ be a homogeneous
coefficient-degree-zero $(1,1)$ symbol with a supplied
compatible ambient finite-PS presentation.  This applies
in particular to LE's actual leading energy $Z_0$.

The smooth test transpose MT2 also agrees on this regular
PS pairing.  Indeed its cocoefficient Radon measures are
the proper transposes of the corresponding original
coefficient measures, by polarization of MT2.  Finite
Borel multiplication by the PS coefficients, with their
actual face values, is precisely the retained regular-PS
module action.  The normal-wall zero-coefficient conclusion
of PT3 removes any ambient representative ambiguity.  Thus
no higher-degree delta action is being inferred from the
smooth transpose alone.

\begin{theorem}[Every critical tangent has the same relevant paired value]
\label{thm:r13-leading-tangent-pairing}
For the cutoff $\rho$ of LE7,
\[
 \begin{split}
 \lim_j\lambda_{\delta_j,S}(E)
       &=-\langle T_0,E\wedge D\rho\rangle,\\
 c_S(P)&=-\langle T_0,Z_0\wedge D\rho\rangle.
 \end{split}
 \tag{PT11}
\]
The second equality holds for every critical tangent of
the protected background.  It is a full tensor pairing,
with all mixed entries and wall traces retained.
\end{theorem}

\begin{proof}
Homogeneity of the coefficient functions of $E$ gives the
exact scaling identity
\[
 \langle T_\delta,E\wedge D\rho\rangle
       =\langle T,E\wedge D[\rho(x/\delta)]\rangle.
 \tag{PT12}
\]
Both $(1,1)$ factors together contribute $\delta^{-4}$
under real pullback, and MT's prefactor $\delta^2$
leaves precisely the $\delta^{-2}$ of the displayed
Hessian.  The original Borel module identification just
explained makes the right side equal to
$\langle M_S(E),D\chi_\delta\rangle$.

MT supplies the uniform trace bounds of PT9 on each
compact real set.  The fixed test $E\wedge D\rho$ is
compatible finite PS.  PT9 therefore applies to the left
side of PT12, even if the tangent current charges walls.
This proves the first equality in PT11.  Taking $E=Z_0$
and using the unconditional leading reduction LE9 proves
the second.  Tangent uniqueness is not asserted or needed.
\end{proof}

If two supplied ambient extensions represent the same actual
symbol, the original Borel module pairings coincide at every
$\delta$.  PT9 consequently gives the same tangent pairing
for both extensions.  This is the precise independence needed
here; support containment alone would not justify multiplying
an arbitrary discontinuous extension by a tangent current.

\paragraph{Remaining analytic realization gate.}
PT11 does not put a critical tangent into the retained analytic
weak-background class, does not construct an original analytic
background from an arbitrary real cone current, and does not
prove its pairing is zero for an actual-zero leading energy.
The realized tangent family is constrained by all original
refinements.  It is not identified with every positive closed
real current on a selected fan or with the full old Killing
dual.  The remaining cancellation is exactly the vanishing
of the paired value in PT11 for those realized backgrounds.

\paragraph{Inputs.}
The real coefficient inequalities and critical compactness are
MT2--MT15, with the middle-positive normalization of PT1.
The original-to-real identification uses the retained R6
regular-PS Borel module, GP's complete proper fibers, and
LE's actual scalar first-jet construction.  The finite-wall
proof and nodal extension are supplied above.  No general
current restriction, arbitrary analytic pushforward, or GK
test extension is used.
